% \section{Symbolic Execution Tools}

\subsection{Owi}\label{section:owi}

Reference \cite{Andrès2024} describes Owi as a toolkit that aggregates functionality to assist WebAssembly developers, one of which enables the compilation of C code to \glsxtrshort{wasm} and running their symbolic interpreter on top of it.

\glsxtrlong{wasm} is a portable binary instruction format for executable programs, designed to be a compilation target for high-level languages like C/C++ and Rust, enabling deployment on the web for client and server applications. WebAssembly is designed to be fast, efficient, and secure, making it an ideal choice for running complex applications in web browsers and other environments \cite{wasm_home,wasm_use_cases}.

% \subsubsection{Installation}

% To properly make use of the tool, a working installation is required and Owi is officially available in two ways. The first and more direct approach is by using the opam, the OCaml Package Manager. The second alternative is through a manual installation using $dune$, the OCaml build system, which can also be installed using $opam$ or manually using $make$ commands.

% The first method resulted in an unusable binary (when it comes to symbolic execution), without any helpful information or help page when running it, therefore is recommendable the use of the second method and building the tool from the repository yourself. To do so, all that needs to be done is running the instructions in \Cref{lst:owi_install} to install \lstinline+opam+, the required dependencies, which include \lstinline+dune+, and finally building and installing \lstinline+owi+.

% \begin{lstlisting}[
% boxpos={ht},
% caption={Owi installation},
% label={lst:owi_install},
% language=bash
% ]
% # install opam 
% apt install opam
% # install owi dependencies
% opam install . --deps-only
% # build owi
% dune build -p owi @install
% # install owi
% dune install
% \end{lstlisting}

\subsubsection{Features}

Owi being a toolchain designed for the development of \glsxtrshort{wasm}, it provides a set of subcommands as listed in \Cref{tab:owi_features}; however, to be able to use the symbolic interpreter packaged with it, we need to also install a restraint solver. Owi supports any Smt.ml supported solver, but the default solver is Z3.

Owi also supports test replayability from an output file generated by the $sym$ subcommand as you can see in \Cref{lst:owi_replay}.

\begin{lstlisting}[
boxpos={ht},
caption={Owi replayability model},
label={lst:owi_replay},
language=bash
]
owi sym main.wat > main.model
owi replay --replay-file=main.model main.wat
\end{lstlisting}

\begin{table}[ht]
    \centering
    \begin{tabular}{r|l}
    \multicolumn{1}{c|}{Subcommand} & \multicolumn{1}{c}{Description} \\
    \hline
    \hline
    $c$ & Compiles C source code to \glsxtrshort{wasm} and runs the \\&\glsxtrshort{wasm} symbolic interpreter on the end result. \\
    $conc$ & Takes \glsxtrshort{wasm} input and runs the concolic interpreter \\&on it.\\
    $fmt$ & Formats \glsxtrlong{wasm} text format files ($.wat$ or \\&$.wast$).\\
    $opt$ & Optimizes \glsxtrlong{wasm} modules.\\
    $run$ & Runs the concrete interpreter.\\
    $replay$ & Replay a \glsxtrshort{wasm} module containing symbols with \\&concrete values in a replay file containing a model.\\
    $script$ & Runs a reference test suite script.\\
    $sym$ & Takes \glsxtrshort{wasm} input and runs the symbolic interpreter \\&on it.\\
    $validate$ & Validates \glsxtrlong{wasm} modules.\\
    $wasm2wat$ & Generate a text format file ($.wat$) from a binary \\&format file ($.wasm$).\\
    $wat2wasm$ & Generate a binary format file ($.wasm$) from a text \\&format file ($.wat$).\\
    \end{tabular}
    \caption{Owi subcommands}
    \label{tab:owi_features}
\end{table}

\subsubsection{Instrumentation}\label{subsection:owi_instrumentation}

The Owi repository maintains a folder with examples of how to use the different subcommands; similarly to KLEE, we have available multiple examples on how to use \lstinline$owi c$ that are detailed in appendix \ref{appendix:owi}.

Looking at the first example, some differences from KLEE are evident, specifically how the tool is informed about which variables are symbolic. KLEE uses the variable address and the type is disregarded in the whole process (the generated test cases detail values for each possible primitive type); Owi on the other hand expects to be told the symbolic variable type, through the available functions:

% \parbox{\textwidth}{
\begin{itemize}
    \item \lstinline$owi_i32()$ - to define 32-bit integer values, with 8-bit and 64-bit variants;
    \item \lstinline$owi_f32()$ - to define 32-bit floating-point values, with a 64-bit variant;
    \item \lstinline$owi_char()$ - to define character values;
    \item \lstinline$owi_bool()$ - to define boolean values.
\end{itemize}
% }

Unlike KLEE functions, these do not support any argument that allows defining a custom name for the symbolic variable; the symbols are identified by order of creation - if we define all positions of an array as symbolic, each position symbolic identifier would correspond to their array position, but if any other symbol is defined before the identifiers would shift by the number of predefined variables.

To run Owi on a simple sample we must include the Owi library in the source code so it is possible to find Owi function implementations and generate a valid WebAssembly module. As visible on \Cref{lst:owi_instrumentation}, the variable \lstinline$x$ of type \lstinline$int$ is instrumented using the routine \lstinline$owi_i32$ (\cref{line:owi_var_x}) and since this program is trying to find the roots of a polynomial, we want the symbolic values that result in \lstinline$poly == 0$, therefore \cref{line:owi_assert_poly} will catch a `problem' every time \lstinline$poly$ equals 0 since the passed condition evaluates to \lstinline$false$.

There are in fact solutions to the defined polynomial and Owi outputs them as visible in \Cref{lst:owi_instrumentation_output}. The output represents a test case, and to replicate the test case one or more of the passed inputs must have a specific set of values. On the example at hand, only one symbol was defined and its state is displayed on \cref{line:owi_symbol_0}. The second value on this line is the symbols' name (\lstinline$symbol_0$), then follows the symbol's type and since an integer was defined its type \lstinline$i32$. Finally, its value is also presented, in this case we can state that 4 is one solution to the polynomial written on \Cref{lst:owi_instrumentation}.

Owi also produces a \lstinline$owi-out/test-suite$ folder and inside it a \lstinline$testcase-N.xml$, where \lstinline$N$ is an index starting on 1, that describes the input value for a symbolic variable for that test case, and a \lstinline$metadata.xml$ that details the arguments of the owi executation such as the source code language, which Owi subcommand produced the file, the input program file (`\textit{poly.c}') and other details on the execution. The generated test cases are relevant to found issues, Owi does not generate test cases on each traced path like KLEE does.

% \noindent\begin{minipage}{\linewidth}
\lstinputlisting[
boxpos=ht,
caption={Polynomial example 1 source code with Owi instrumentation},
label={lst:owi_instrumentation},
escapechar=|,
numbers=left,
escapechar=|,
breaklines=true,
% upquote=true,
showstringspaces=false,
postbreak=\mbox{\textcolor{red}{$\hookrightarrow$}\space},
]{code-snippets/chap3/owi_instrumentation}
% \end{minipage}

\lstinputlisting[
boxpos=ht,
caption={Polynomial example 1 output},
label={lst:owi_instrumentation_output},
escapechar=|,
numbers=left,
escapechar=|,
breaklines=true,
% upquote=true,
showstringspaces=false,
postbreak=\mbox{\textcolor{red}{$\hookrightarrow$}\space},
]{code-snippets/chap3/owi_instrumentation_output}

\subsubsection{Running Owi}

Similarly to KLEE, Owi provides a set of methods to inform the symbolic interpreter in an attempt to improve the results, examples of this are \lstinline{owi_assume()} and \lstinline{owi_i32()} for int32 variables (and others for the respective primitive types such as \lstinline{owi_char()}).

Once the sources are ready, running the symbolic interpreter is very simple with Owi by running \autoref{lst:owi-command}. Since this is a toolkit it contains many subcommands, the most relevant for this study being the `$c$' subcommand that allows running C code directly.

\begin{lstlisting}[
caption={Running owi on a program source.},
label={lst:owi-command}
]
owi c main.c
\end{lstlisting}

However, if your program requires any sort of custom option or if you just want to experiment with them, we can `ask' Owi how they compile the sources. To do that, all one needs to do is run \lstinline$owi c main.c --debug$ on invalid C files and will find in the tool output the command in \Cref{lst:owi_clang}.

% \noindent\begin{minipage}{\linewidth}
\begin{lstlisting}[
boxpos=ht,
caption={Owi abstracted compilation command (clang).},
label={lst:owi_clang},
breaklines=true,
escapechar=|,
]
clang -O3 --target=wasm32 -m32 -ffreestanding \
    |\mbox{\textcolor{red}{$\hookrightarrow$}\space}|--no-standard-libraries -Wno-everything -flto=thin \
    |\mbox{\textcolor{red}{$\hookrightarrow$}\space}|-Wl,--entry=main -Wl,--export=main -Wl,--lto-O0 \
    |\mbox{\textcolor{red}{$\hookrightarrow$}\space}|-Wl,-z,stack-size=8388608 \
    |\mbox{\textcolor{red}{$\hookrightarrow$}\space}|-I $HOME/.opam/default/share/owi/libc -o main.wasm \ 
    |\mbox{\textcolor{red}{$\hookrightarrow$}\space}|$HOME/.opam/default/share/owi/binc/libc.wasm \
    |\mbox{\textcolor{red}{$\hookrightarrow$}\space}|main.c
\end{lstlisting}
% \end{minipage}

Owi also integrates Frama-C's E-ACSL plugin. This plugin enables the detection of implementation issues, or even custom heuristics for complex bugs, through the definition of a set of rules that can create constraints on acceptable inputs (`requires') and outputs (`ensures'). The \lstinline$owi c$ subcommand and to enable all one needs to do is pass the argument \lstinline$--e-acsl$. To see how Frama-C is being executed the process is similar as to the compilation, but now with the E-ACSL option, and it looks like \Cref{lst:owi_frama_c}. This plugin enables the detection of implementation issues, or even custom heuristics for complex bugs, through the definition of a set of rules that can create constraints on acceptable inputs (`requires') and outputs (`ensures').

\begin{lstlisting}[
boxpos=ht,
caption={Owi abstracted Frama-C instrumentation of E-ACSL.},
label={lst:owi_frama_c},
breaklines=true,
escapechar=|,
% upquote=true,
showstringspaces=false,
]
frama-c -e-acsl -no-frama-c-stdlib \
    |\mbox{\textcolor{red}{$\hookrightarrow$}\space}|-kernel-warn-key CERT:MSC:38=inactive -verbose 2 \
    |\mbox{\textcolor{red}{$\hookrightarrow$}\space}|-cpp-extra-args="-I$HOME/.opam/default/share/owi/libc" \
    |\mbox{\textcolor{red}{$\hookrightarrow$}\space}|main.c -then-last -print -ocode main.instrumented.c
\end{lstlisting}

After compiling either \lstinline$main.c$ or \lstinline$main.instrumented.c$, the resulting \lstinline$.wasm$ object can be tested with Owi using the command on \Cref{lst:owi-sym-command}.

\begin{lstlisting}[
caption={Running owi on a compiled program.},
label={lst:owi-sym-command},
]
owi sym main.wasm
\end{lstlisting}

Executing either \lstinline$c$ or \lstinline$sym$ subcommands generates some files inside the folder `owi-out' created on the current working directory: \lstinline$owi sym$ only generates test case files named `testcase-n.xml' being $n$ the number of results detected and starting from 1. The subcommand \lstinline$owi c$ also generates a file describing the test cases that fail execution, but also outputs an additional `metadata.xml' 
% as described in \Cref{tab:owi_c_owi_out}
, and a `a.out.wasm' file (the compiled sources); however, the test case and metadata files are generated in a subfolder called `test-suite' while the \glsxtrlong{wasm} module is created on the current working directory (at the same level as `owi-out' folder).
For a visual reference of this file structure take a look at \Cref{lst:owi-generated-files}, consider `main.c' the input sources and `main.wasm' pre-compiled before running Owi.

% \begin{table}[ht]
%     \centering
%     \begin{tabular}{r|c|l}
%         \multicolumn{1}{c|}{Field} & \multicolumn{1}{c|}{Example} &\multicolumn{1}{c}{Description} \\
%         \hline
%         \hline
%         \lstinline$sourcecodelang$ & `C' & The programming language in \\&&which the input sources are \\&&written \\
%         \lstinline$producer$ & `owic' & The subcommand responsible for \\&&this metadata file.\\
%         \lstinline$specification$ & N/A & \multicolumn{1}{c}{No information on this}\\&& \multicolumn{1}{c}{field was found}\\
%         \lstinline$programfile$ & `main.c' & The input source files\\
%         \lstinline$programhash$ & b663247... & The compiled wasm module hash\\
%         \lstinline$entryfunction$ & `main' & \glsxtrshort{wasm} modules requires by \\&&default a specific entry point, \\&&in this case Owi defines `main'\\
%         \lstinline$architecture$ & `32bit' & The target architecture (defined \\&&by Owi)\\
%         \lstinline$creationtime$ & timestamp & The day and time this metadata \\&&file was created\\
%     \end{tabular}
%     \caption{Owi `metadata.xml' file contents.}
%     \label{tab:owi_c_owi_out}
% \end{table}

\noindent
% \begin{minipage}{\textwidth}
\begin{lstlisting}[
caption={Owi generated files structure.},
label={lst:owi-generated-files},
literate={├}{{\textSFviii}}1 {─}{{\textSFx}}1 {└}{{\textSFii}}1 {│}{{\textSFxi}}1,
]
# owi c ...
.
├── a.out.wasm
├── main.c
└── owi-out
   └── test-suite
      ├── metadata.xml
      └── testcase-1.xml

# owi sym ...
.
├── main.c
├── main.wasm 
└── owi-out
   └── testcase-1.xml
\end{lstlisting}
% \end{minipage}
